%% kapa: exam=1 round=29 year=2018 period=1 form=A subject=civil_law subject_ko=민법 seq=39 number=39
%% source: Q-Net 감정평가사 기출문제 (cst00309) artlSeq=5202808 file=extracted/1008/5202808/2018년 제29회 감정평가사 1차시험 기출문제/2018년 제29회 감정평가사 1차시험 기출문제/1교시 A형.pdf
%% conversion: kordoc 4.24.51 -> tools/gen1.py (auto); review=proofread-sonnet; flags=none
\begin{question}{39}
법정지상권에 관한 설명으로 옳은 것은? (다툼이 있으면 판례에 따름)
\begin{choices}
  \item 관습상 법정지상권이 성립하려면 토지와 그 지상건물이 원시적으로 동일인의 소유에 속하고 있어야 한다.
  \item 토지에 저당권이 설정될 때에 그 지상건물이 미등기인 경우, 저당권 실행으로 토지와 건물의 소유자가 상이하게 되더라도 법정지상권은 인정될 수 없다.
  \item 환지처분으로 인하여 토지와 그 지상건물의 소유자가 달라진 경우에도 관습상 법정지상권은 인정될 수 있다.
  \item 토지와 그 지상건물의 소유자가 달라질 때, 토지의 사용에 대하여 당사자 사이에 특약이 있는 경우, 관습상 법정지상권은 인정될 수 없다.
  \item 나대지상에 채권담보를 위한 가등기가 경료된 후에 대지소유자가 그 지상에 건물을 신축하였고, 그 후에 가등기에 기한 본등기가 경료되어 대지와 건물의 소유자가 달라진 경우 관습상 법정지상권이 성립될 수 있다.
\end{choices}
\end{question}
\answer{4}
